%! Characteristics of Quadratic Functions — Unit 2, Lesson 2.0 — Lesson Plan
% GRADUAL-RELEASE plan (warm-up / guided notes & practice / debrief / close & start homework) on
% the Main Ideas / Notes table shape (LESSON_SHAPE.md §1, §5). NO GROUP ACTIVITY, NO EXIT TICKET,
% NO INDEPENDENT PRACTICE SET IN THE NOTES — the homework is the individual practice, and the
% period ends with students starting it. No tiers, no extension box. Teacher-facing, 10pt; every
% student component is 12pt. Timing 5 / 35 / 10 / 10.
% THIS IS A LESSON 0 — a characteristics lesson. Pure graph READING: no solving, no factoring.
% 2.1 graphs from the equation; 2.2 factors.
% Prose trimmed 2026-09-17 (first build came out 7pp); all §5 sections and all three teacher
% notes survive — only wording was cut.
\documentclass[10pt]{article}
\usepackage{algebra2-article}
\usepackage{algebra2-boxes}
\usepackage{graphicx}   % -article does NOT load graphicx; needed for thumbnails/screenshots
\graphicspath{{images/}}
\newcommand{\TallMath}[1]{$\displaystyle #1\rule[-1.4em]{0pt}{3.2em}$}
\newcommand{\CourseName}{Algebra 2}
\newcommand{\MeetingLength}{60 minutes}
\newcommand{\UnitNumberName}{Unit 2: Quadratic Functions \quad}
\newcommand{\LessonNumberName}{Lesson 2.0: Characteristics of Quadratic Functions}

\begin{document}
\begin{center}
    {\Large\bfseries \CourseName \\
    {\small\bfseries \UnitNumberName \LessonNumberName}}
\end{center}

\begin{tcolorbox}[colback=forestbg, colframe=forest, boxrule=0.9pt, arc=2mm,
    left=2mm, right=2mm, top=1.5mm, bottom=1.5mm]
    \textbf{Primary Objective:} Students read \emph{every} characteristic of a quadratic function
    off a \textbf{pre-drawn graph} or a \textbf{table} --- \textbf{vertex} (turning point),
    \textbf{maximum} or \textbf{minimum value}, \textbf{axis of symmetry}, \textbf{zeros}
    ($x$-intercepts), \textbf{$y$-intercept}, \textbf{increasing} and \textbf{decreasing}
    intervals, \textbf{domain} and \textbf{range}, and \textbf{end behavior} --- and justify each
    with \textbf{symmetry}. \emph{Nothing is solved and nothing is factored today}: this is the
    Unit~2 Lesson~0, the read-the-graph toolkit for the parabola.
    \\[2pt]
    \textbf{Standards (2023 VA SOL):} \textbf{A2.F.2a, A2.F.2c, A2.F.2d, A2.F.2f, A2.F.2g} ---
    the characteristics of \emph{quadratic} functions, algebraically and graphically: domain,
    range, zeros, intercepts (\textbf{a}); intervals where the function is increasing or
    decreasing (\textbf{c}); \textbf{absolute maximum and absolute minimum} (\textbf{d}); for any
    $x$ in the domain, determine $f(x)$ from a graph or equation and explain its meaning in
    context (\textbf{f}); and \textbf{end behavior} (\textbf{g}). \textbf{a}, \textbf{c} and
    \textbf{f} are revisited from Unit~1 on a new graph; \textbf{d} and \textbf{g} are
    \emph{new} --- the parabola is the first family in this course that requires them.
    \\[2pt]
    \textbf{Lesson model:} \emph{Gradual release --- I do, we do, you do --- all inside the guided
    notes.} There is no group activity. The notes name the vocabulary as they teach it and deliver
    the instruction in \textbf{one Main Ideas / Notes table} on one context (a ball off a 48-foot
    ledge): four instruction rows --- \emph{The turning Point}, \emph{The mirror Line},
    \emph{Reading the Whole Graph}, \emph{How high vs.\ Where} --- each a definition, a pre-drawn
    display, and a grid of problems (1--9) where you work the first and the class works the rest;
    then the Guided Practice row (10--13), on a graph with \emph{no story}, is worked together
    with students holding the pen. \textbf{The homework is the individual practice}, started in
    class while you circulate. The debrief is \emph{spoken}; there is no practice set and no
    objective box on the notes page --- the cover carries the targets.
\end{tcolorbox}

\renewcommand{\arraystretch}{1.4}
\boxguard
\begin{skillbox}[Priority Ideas \& Skills]{goldbox}
\begin{tabularx}{\linewidth}{@{}X|X@{}}
\textbf{Priority Ideas \& Skills} & \textbf{Key Understandings (the ``why'')} \\[2pt]
\hline
\vspace{-6pt}
\begin{itemize}[leftmargin=1.1em, nosep, after=\vspace{2pt}]
  \item Name the \textbf{vertex} and say whether it is a maximum or a minimum point, from the
        sign of $a$.
  \item State the \textbf{maximum} or \textbf{minimum value} --- with units --- \emph{and} the
        input at which it occurs, as two different answers.
  \item Write the \textbf{axis of symmetry} $x=h$; use symmetry to pair equal outputs and to
        locate the vertex from a table or from the two zeros.
  \item Read the \textbf{zeros} (two, one, or none) and the \textbf{$y$-intercept}.
  \item Name the \textbf{increasing} and \textbf{decreasing} intervals, split at the axis; give
        \textbf{domain} (all reals) and \textbf{range} ($y\ge k$ or $y\le k$).
  \item Describe \textbf{end behavior}: $a>0$ both ends rise, $a<0$ both ends fall.
\end{itemize}
&
\vspace{-6pt}
\begin{itemize}[leftmargin=1.1em, nosep, after=\vspace{2pt}]
  \item A parabola is completely described by its turning point and which way it opens; every
        other feature is read off that, and symmetry is the reason.
  \item \textbf{Target misconception:} that the maximum or minimum \emph{value} is the
        $x$-coordinate of the vertex --- ``the ball's maximum is $1$.'' \textbf{The value is an
        \emph{output} (a $y$-value); \emph{where} it happens is an \emph{input}.} Both numbers
        answer \emph{different} questions: the vertex $(1,64)$ says the greatest height is
        $64$~ft, attained at $t=1$~s. Notes~8 is the classmate claim, 9 pairs value against
        where, GP~13 tests it on new ground.
  \item The axis is where the function \emph{turns}, so it is exactly the boundary between the
        increasing and the decreasing interval --- one fact, two features.
  \item A parabola's two ends always go the \emph{same} way; a line's go opposite ways. First
        end-behavior contrast in the course.
\end{itemize}
\end{tabularx}
\end{skillbox}

\boxguard
\begin{skillbox}[{Vocabulary, Concepts, \& Theorems --- taught directly in the guided notes}]{sky}
\renewcommand{\arraystretch}{1.3}
\begin{tabularx}{\linewidth}{@{}l X@{}}
\textbf{Parabola} & The graph of $f(x)=ax^2+bx+c$, $a\ne0$; opens up when $a>0$, down when $a<0$. \\
\textbf{Vertex (turning point)} & Where the parabola turns: lowest point if it opens up, highest if it opens down. \\
\textbf{Maximum / minimum value} & The \emph{output} at the vertex --- a $y$-value with units; the $x$-coordinate is only \emph{where} it occurs. \\
\textbf{Axis of symmetry} & The line $x=h$ through the vertex; inputs equally far either side have equal outputs. \\
\textbf{End behavior} & What the outputs do far left and far right: $a>0$ both ends rise, $a<0$ both ends fall --- always the same way. \\
\textbf{Zeros / $x$-intercepts} & \emph{(Revision from Unit 1.)} Inputs where $f(x)=0$; two, one, or none, symmetric about the axis. \\
\end{tabularx}
\end{skillbox}

\begin{fixedskillbox}[Lesson at a Glance --- \MeetingLength]{forestbg}
\renewcommand{\arraystretch}{1.25}
\begin{tabularx}{\linewidth}{@{}l c X X@{}}
\textbf{Phase} & \textbf{Min} & \textbf{Students} & \textbf{Teacher} \\
\hline
Warm-Up & 5 & Three prior skills, alone, then check with a neighbour & Circulate; debrief item~3
  aloud (``the highest point is $(2,4)$ --- is the \emph{highest} answer the $x$- or the
  $y$-value?'') and leave it hanging \\
Guided Notes \& Practice & 35 & Fill the vocabulary box and the table row by row as each term is
  named, working the rest of each grid themselves; then Guided Practice (10--13) together, pen in
  their hand & \textbf{I do} the hook, then each row's definition, display and first problem
  (1, 3, 5, 8) (20) $\to$ \textbf{we do} the rest of each grid, then the Guided Practice row (15)
  on $g(x)=-(x+2)^2+9$ \\
Debrief & 10 & Pens down, papers up --- answer aloud, nothing to fill in & Put $(1,64)$ back on
  the board; take the ``the maximum is $1$'' reflex and the ``domain is only where the ball
  flies'' error \\
Close \& Start Homework & 10 & \textbf{Start the homework in class}, alone and silent & Launch
  item~1 aloud, then circulate for your formative read; preview Lesson~2.1; the rest is due the
  first class after two study halls \\
\end{tabularx}
\end{fixedskillbox}

\begin{fixedskillbox}[Warm-Up --- Activate Prior Knowledge (5 min)]{forestbg}
\begin{minipage}[t]{0.48\linewidth}
    \textbf{The three items and what each seeds}
    \begin{itemize}
        \item \textbf{1.} Evaluate $f(x)=x^2-2x-3$ at $-1,0,1,3$. \textbf{Seeds:} substitution
              with negatives and squares --- and it quietly hands them the anchor curve's zeros
              ($-1$, $3$) and vertex output ($-4$) before they see the graph.
        \item \textbf{2.} Read a \emph{linear} graph from Unit~1: intercepts, increasing or
              decreasing, domain. \textbf{Seeds:} ``read the graph, name the feature'' is already
              the move; only the shape changes today.
        \item \textbf{3.} A pre-drawn parabola with its high point circled: increasing until
              \underline{\hspace{0.5cm}}, decreasing after \underline{\hspace{0.5cm}}, highest
              point at \underline{\hspace{0.5cm}} --- \emph{and is the ``highest'' answer the
              $x$-value, the $y$-value, or both?} \textbf{Seeds:} the whole lesson. Row~1 answers
              the first two; row~4 settles the third.
    \end{itemize}
\end{minipage}
\hfill
\begin{minipage}[t]{0.48\linewidth}
    \textbf{Running it}
    \begin{itemize}
        \item \textbf{Debrief item~3 aloud, briefly}, and land only this: \emph{the point is
              $(2,4)$ --- hold on to which of those two numbers ``highest'' is asking for.} Do not
              answer it; the hook's vote is worth more than the ten seconds.
        \item Item~1's $f(1)=-4$ is the one to point back at in row~2: it is the anchor curve's
              vertex output, found before the picture.
        \item If it runs long, cut item~2's domain question and keep item~3 whole.
        \item The page is 12pt like the rest of the packet and fits one side. If you add to it,
              check that it still does.
    \end{itemize}
\end{minipage}
\end{fixedskillbox}

\boxguard
\begin{skillbox}[Guided Notes \& Practice --- I do, then we do (35 min)]{forestbg}
\textbf{This is the whole lesson.} No group activity, and no independent practice set on this
page: \textbf{the homework is the individual practice}, started in class.
\begin{multicols}{2}
\textbf{Hook (3 min, page 1).} ``A ball is launched from a 48-foot ledge; its height after $t$
seconds is $h(t)=-16t^2+32t+48$.'' Read the pre-drawn graph together: starts at $48$~ft, turns at
$(1,64)$, lands at $t=3$. \textbf{Say out loud that the $h$-axis is scaled --- one grid unit is
$16$ feet} --- or students read the vertex as $(1,4)$. Then the vote: \textbf{``The ball's maximum
is: circle one --- $1$ or $64$ --- and what are the units?''} Hand count, one reason each,
\emph{do not settle it}. Problem~8 settles it; 13 tests it again. Fill the vocabulary rows
\emph{as you go}, never in advance.

\textbf{I do, row by row (about 17 min).} For each row: read the definition aloud, mark up the
display, work the \emph{first} problem of the grid, hand the rest to the room.

\emph{The turning Point (5 min).} Parabola; the sign of $a$; the vertex as the turning point ---
here a \emph{maximum} point because $a=-16<0$. Work \textbf{1}; the class works \textbf{2}
($(1,64)$, a maximum point).

\emph{The mirror Line (4 min).} The axis $x=h$; equal outputs either side. Switch to the unit
anchor $f(x)=x^2-2x-3=(x-1)^2-4$ --- vertex $(1,-4)$, axis $x=1$, zeros $-1$ and $3$,
$y$-intercept $(0,-3)$ --- \textbf{and say it returns in 2.1--2.5}. Fill the mirror-pair table
\emph{across} the axis. Work \textbf{3} ($f(2)=-3$ because $f(0)=-3$); the class works \textbf{4}
(average the zeros: $1$).

\emph{Reading the Whole Graph (4 min).} Everything else off the same annotated anchor: zeros (two,
one, or none), $y$-intercept, increasing/decreasing split \emph{at} the axis, domain all reals,
range $y\ge k$ or $y\le k$. Work \textbf{5}; the class works \textbf{6} and the feature card
\textbf{7}.

\emph{How high vs.\ Where (4 min) is the lesson.} Two sentences: \emph{the value is an output, the
location is an input.} \textbf{8 is the crux:} let someone defend ``the maximum is $1$'' and let
the room catch it --- $1$ answers \emph{when}, in seconds; the value is $64$, in feet. Ask for the
\emph{units}, do not just repeat the rule. Then end behavior off the two small parabolas. The
class works \textbf{9}.

\columnbreak
\textbf{We do (about 15 min).} Guided Practice on $g(x)=-(x+2)^2+9$, pre-drawn, no story: vertex
$(-2,9)$, opens down, zeros $-5$ and $1$, $y$-intercept $(0,5)$. \textbf{10:} opens which way,
vertex, axis. \textbf{11:} zeros and $y$-intercept, each zero three units from $x=-2$.
\textbf{12:} increasing $(-\infty,-2]$, decreasing $[-2,\infty)$, domain all reals, range
$y\le9$, both ends fall. \textbf{13 is the crux transferred and the problem to protect:} the
maximum \emph{value} is $9$, occurring \emph{at} $x=-2$; the classmate answering ``$-2$'' gave the
\emph{input} instead of the \emph{output}.

\textbf{The pen must actually change hands.} This is the only place today you can watch a student
decide, so do not narrate it. Put 10's axis on them cold, and take a hand count on 13's ``$9$ or
$-2$'' before anyone speaks.

Circulate and demand the reason, not the word:
\begin{itemize}[leftmargin=1.1em, nosep]
  \item \textbf{11:} ``How far is $-5$ from the axis? How far is $1$? Why must they match?''
  \item \textbf{12:} ``Where does the curve turn? So where does increasing stop?''
  \item \textbf{13, the crux:} ``Is $9$ an input or an output? Which question does $-2$ answer?''
\end{itemize}
Pick two papers for the debrief --- one 13 that wrote \emph{value $9$, at $x=-2$}, one that wrote
\emph{the maximum is $-2$}. \textbf{Your formative read comes twice}: here, and again in the last
ten minutes.
\end{multicols}
\end{skillbox}

\boxguard
\begin{skillbox}[Debrief --- whole class, spoken (10 min)]{forestbg}
Pens down, papers up. \textbf{This debrief is spoken} --- nothing at the end of the notes to fill
in.
\begin{multicols}{2}
\begin{enumerate}[leftmargin=1.3em, nosep, itemsep=3pt]
  \item \textbf{Put $(1,64)$ back up} and make a student say both sentences: \emph{the maximum
        value is $64$ feet}, and \emph{it happens at $t=1$ second.} Settle the hook's vote by
        naming the units. That is the lesson in one line.
  \item Put up the two problem~13s. ``$-2$'' gave \emph{where}, not \emph{how high}; make the room
        say which axis each number lives on. Point back at $(1,64)$ --- same mistake, new curve.
  \item \textbf{Symmetry does the work.} Why must $-5$ and $1$ be the same distance from $x=-2$?
        Why does averaging the zeros find the axis? Then $f(0)=f(2)=-3$ on the anchor.
  \item The domain reflex last: the \emph{function} has domain all real numbers; only the
        \emph{ball} stops at $t=3$. Close on end behavior --- both ends the same way, always.
\end{enumerate}
\columnbreak
\textbf{Two things to demand aloud.} First, \textbf{the vertex sentence without the notes} ---
``$(1,64)$ means the greatest height is $64$ feet, one second after launch.'' Second,
\textbf{what the axis does} --- it is the mirror \emph{and} the place the curve turns --- because
homework~3 finds a vertex from a table on nothing but that.

\textbf{If short}, cut item~4 then item~3; item~2 cannot be cut. \textbf{Do not let the debrief
eat the last ten minutes} --- the homework start is the individual practice, not padding.
\end{multicols}
\end{skillbox}

\boxguard
\begin{skillbox}[Homework --- scored, started in class, due the first class after two study halls]{goldbox}
Six items spanning A2.F.2a/c/d/f/g, two pages, no extension: \textbf{1} the core read off one
pre-drawn parabola (vertex, axis, opens which way, minimum or maximum \emph{and its value}, zeros,
$y$-intercept, increasing/decreasing, domain, range); \textbf{2} the \emph{contrast pair}, one
parabola opening up beside one opening down --- which value each has, what it \emph{is} and
\emph{where} --- closing with a ``why?''; \textbf{3} the same skill \emph{off a table} --- constant
\textbf{second differences} confirm it is quadratic, then the table's symmetry gives the axis and
vertex; \textbf{4} the special case and its boundary --- a vertex \emph{on} the $x$-axis (one
zero), a parabola that never reaches it (none), and ``why can the answer never be three?'';
\textbf{5} a punt, $h(t)=-16t^2+48t+4$, with a \texttt{work} block for $h(1.5)$, an
interpret-the-answer follow-up, and a domain / end-behavior sanity check; \textbf{6} an SOL-style
multiple-choice item on maximum \emph{value} vs.\ \emph{where} --- the formative check.

\textbf{Start it in the last ten minutes, in class and alone.} Item~1's ``value'' blank and
item~2's ``why?'' are where the misconception shows. \textbf{Item~6's sort:} (A) $3$ is the
\emph{input} (the target misconception), (C) $6$ is the far zero, (D) $96$ is a coefficient.

\textbf{Two pages is a ceiling at 12pt}, so a seventh item was cut rather than spilled --- the
parabola-vs-line end-behavior comparison lives in the debrief and in notes~9 instead.
\textbf{DeltaMath override:} reading vertex, axis, intercepts and intervals off a quadratic graph
is \emph{well} covered there --- swap in a set if you prefer auto-graded practice; that replaces
this page and strikes its score cell on the cover. Never assign both; items~3 and~4 are thinner
there, so keep them if you do. \textbf{Preview:} Lesson~2.1, \emph{Graphing Quadratic Functions}
--- the same curve in three costumes, and the graph built \emph{from} the equation.
\end{skillbox}

\boxguard
\begin{skillbox}[Watch For (while circulating)]{redbox}
\begin{itemize}[leftmargin=1.2em, nosep]
  \item \textbf{Giving the $x$-coordinate as the maximum or minimum \emph{value}} (notes~8, 9,
        GP~13, HW~1, 2, 6) --- the target misconception. Probe: \emph{``Input or output? Which
        question does the other number answer?''} In a context, demand the units.
  \item \textbf{Writing the vertex as a single number} instead of a point (notes~2, GP~10).
        Probe: \emph{``A point needs two coordinates --- how high, and when?''}
  \item Saying the domain is only the drawn part of the graph (notes~7, GP~12, HW~5). The
        \emph{function's} domain is all reals; the \emph{model} stops when the ball lands.
  \item Naming increasing/decreasing intervals in $y$-values, or splitting somewhere other than
        the axis (notes~6, GP~12). Probe: \emph{``Where does it turn?''}
  \item ``One end up, one end down'' --- borrowed from lines (notes~9, GP~12, HW~2). A parabola's
        ends always go the \emph{same} way.
  \item Expecting every parabola to have two zeros (HW~4). Two, one, or none --- never three,
        because the curve turns only once.
\end{itemize}
Cold-call prompts: ``Input or output?'' ``How far is that from the axis --- and what is on the
other side?'' ``Which way does it open, and how do you know from the picture?'' ``Both ends ---
same way or opposite?''
\end{skillbox}

\boxguard
\begin{skillbox}[Close \& Start Homework (10 min)]{goldbox}
\textbf{Students start the homework here, in class, alone and silent --- this is the individual
practice, and the first minutes are your best formative read.} Hand the launch first --- six items
on paper, scored, due the first class after two study halls (announce the date in class and on
TurtleNet) --- and read item~1's vertex aloud in thirty seconds so nobody stalls. Then circulate on
item~1's ``value'' blank and item~2's ``why?'', sorting into three piles: gave the output with
units and named \emph{where} separately (got it); gave the $x$-coordinate as the value (the target
misconception, still alive); gave a point where a number was asked for or the reverse (a notation
fix, not a concept one). Close on what changed: \emph{``You walked in able to read a line. You walk
out able to read a parabola --- where it turns, how high that is, where it crosses, where it rises
and falls, and what both ends do --- and to say which answers are inputs and which are outputs.''}
\end{skillbox}

% ── Teacher notes — one per component, in packet order (three: no activity, no practice set) ──
\begin{teachernote}[Warm-Up]
Item~1 is the mechanical core --- substitution with a negative input and a square --- and it is
not a throwaway: $f(x)=x^2-2x-3$ at $-1,0,1,3$ gives $0,-3,-4,0$, which is the anchor curve's two
zeros and its vertex output, found before they have seen the picture. Point back at that in the
\emph{mirror Line} row. Item~2 is Unit~1's graph read on a line, so only the shape is new today.
Item~3 is the make-or-break item: the parabola is pre-drawn with its high point circled, the first
two blanks are answered by row~1, and the last question --- \emph{is the ``highest'' answer the
$x$-value, the $y$-value, or both?} --- is the target misconception asked before any vocabulary
exists. Debrief it in one sentence and \textbf{leave it hanging}; the hook's vote is thirty
seconds away. One page at 12pt, blank and keyed, two small pre-drawn figures, no sketching.
\end{teachernote}

\begin{teachernote}[Guided Notes \& Practice]
\textbf{This one document is the lesson --- pace it 3 / 17 / 15, then 10 for the debrief and 10
for the homework start.} Page~1 is the five vocabulary rows and the hook; fill the rows \emph{as
each term is named} and take the ``$1$ or $64$'' vote without resolving it. \textbf{Say out loud
that the projectile graph's $h$-axis is scaled} --- one grid unit is $16$ feet --- because on a
plain unit grid this curve is a page-tall spike, and students who do not hear it read the vertex
as $(1,4)$.

Then the table, row by row: definition, display, first problem (1, 3, 5, 8), rest to the room.
\emph{The turning Point} must move: parabola, the sign of $a$, the vertex, problems 1--2.
\emph{The mirror Line}: the axis $x=1$ on the anchor $f(x)=x^2-2x-3=(x-1)^2-4$, the mirror-pair
table filled \emph{across} the axis, problems 3--4 --- and \textbf{tell them this curve comes
back}, it is the Unit~2 anchor for 2.1--2.5. \emph{Reading the Whole Graph}: zeros (two, one, or
none), $y$-intercept, increasing/decreasing split \emph{at} the axis, domain all reals, range
$y\ge k$ or $y\le k$, and the feature card, problem~7, which sits beside the graph so the answers
come straight off the picture.

\textbf{Spend the time in \emph{How high vs.\ Where}}, which carries the misconception. Two
printed sentences, then \textbf{problem~8}: let someone defend ``the maximum is $1$'' and let the
room catch it. Do not resolve it by repeating the rule --- point at the vertex and ask for the
\emph{units}. Seconds answer \emph{when}; feet answer \emph{how high}. Close the row on the two
small parabolas: $a>0$ both ends rise, $a<0$ both ends fall, and a parabola's ends always agree
--- the first place in the course where a line's behavior is the \emph{contrast} rather than the
model. Problem~9 pairs value against where and asks the end behavior of both displays.

\textbf{Guided Practice (10--13) is where the pen changes hands}, on $g(x)=-(x+2)^2+9$ with no
story: 10 vertex and axis, 11 zeros and the symmetry check, 12 intervals, domain, range and end
behavior, \textbf{13 the crux transferred} --- value $9$, occurring at $x=-2$, and what the
classmate answering ``$-2$'' has confused. If time is short, cut 12, never 13.

\textbf{There is no independent practice set on this page, deliberately --- the homework is the
individual practice}, and the period ends with students starting it in front of you. \textbf{Nothing
here asks a student to sketch a graph}: every display is pre-drawn to read, annotate, or a table to
fill. No \texttt{work} blocks in this component --- every answer space is a fixed-height grid cell,
so blank and key cannot drift.

\textbf{Debrief (10 min), spoken.} $(1,64)$ with its units, then the two problem~13s; item~2 is
the one move that cannot be cut.
\end{teachernote}

\begin{teachernote}[Homework]
\textbf{Scored; due the first class after two study halls --- announce the date in class and on
TurtleNet.} The ten minutes at the end are a \emph{start}: put students on 1 and 2 while you can
still catch a wrong start, and let 3--6 go home. Item~1 is the whole feature list off one graph;
the blank to watch is the \emph{value}, where you will see the $x$-coordinate written in. Item~2's
``why?'' should say \emph{because it opens down, the vertex is the highest point} --- a student who
writes only ``because the vertex is there'' has not used the direction. Item~3 runs off a table:
first differences $6,2,-2,-6$, second differences all $-4$ (constant, so quadratic), then
$y(1)=y(3)=3$ pairs across the axis $x=2$ and the vertex is $(2,5)$; push back if they read the
vertex as the largest $y$ in the list without the symmetry argument. Item~4: $y=(x-2)^2$ touches at
exactly one point, $y=x^2+2$ never reaches the axis, and a parabola opening up with its vertex
\emph{above} the axis has \emph{no} zeros --- three is impossible because the curve turns only
once. Item~5 has the one \texttt{work} block ($h(1.5)=40$), the maximum height $40$~ft at
$t=1.5$~s, and the sanity question: the model stops meaning anything after the ball lands (about
$t\approx3.1$~s) even though the function keeps falling forever. \textbf{Sort item~6 into three
piles:} chose B (got it); chose A (gave the input $3$ --- the target misconception, still alive);
chose C or D (read a zero or a coefficient off the rule). If more than a third land outside B, open
Lesson~2.1 with a one-minute re-look at the vertex as \emph{(where, how high)}. \textbf{DeltaMath
override:} a swapped-in set replaces this page and strikes its score cell on the cover --- never
assign both.
\end{teachernote}

\end{document}
